\answersfor{ch:losses}

\begin{solution}{ex:losses:1}
\textbf{Errors} $e = \hat{x} - x$:
\begin{align*}
\text{A, band 1: } & 0.13 - 0.10 = +0.03\\
\text{A, band 2: } & 0.16 - 0.20 = -0.04\\
\text{B, band 1: } & 0.30 - 0.30 = \phantom{+}0.00\\
\text{B, band 2: } & 0.09 - 0.05 = +0.04
\end{align*}
\textbf{Absolute values:} $0.03, 0.04, 0.00, 0.04$.

\textbf{Band means per pixel:}
\begin{align*}
\text{A: } & (0.03 + 0.04)/2 = 0.07/2 = 0.035\\
\text{B: } & (0.00 + 0.04)/2 = 0.04/2 = 0.020
\end{align*}
\textbf{Mean over the two graded pixels:} $(0.035 + 0.020)/2 = 0.055/2 = 0.0275$.

Check with the second form: the sum of all four absolute errors is $0.03 + 0.04 + 0.00 + 0.04 =
0.11$, and $n = C\abs{M} = 2 \times 2 = 4$, so $0.11/4 = 0.0275$.

\textbf{Gradient.} Every graded entry gets $\sign(e)/n$ with $1/n = 1/4 = 0.25$:
\[
\frac{\partial L_1}{\partial \xhat_A} = (+0.25,\ -0.25), \qquad
\frac{\partial L_1}{\partial \xhat_B} = (0,\ +0.25).
\]
Band 1 of B is exact, so $\sign(0) = 0$ and it gets no push.
\end{solution}

\begin{solution}{ex:losses:2}
\textbf{The angle.} $\ip{\vect{x}}{\xhat} = 1 \times 1 + 0 \times 1 = 1$. $a = \norm{\vect{x}} =
\sqrt{1 + 0} = 1$. $b = \norm{\xhat} = \sqrt{1 + 1} = \sqrt{2} = 1.414214$.
\begin{align*}
u &= 1/(1 \times 1.414214) = 0.707107\\
\theta &= \arccos(0.707107) = \pi/4 = 0.785398 \text{ rad}\\
\theta &= 0.785398 \times 180/\pi = 45\text{\textdegree}
\end{align*}
\textbf{The cosine's gradient}, \cref{eq:losses:du}: $\partial u/\partial\xhat = \frac{1}{b}\big(\vect{x}/a - u\,\xhat/b\big)$.
\begin{align*}
\vect{x}/a &= (1, 0)\\
\xhat/b &= (1/1.414214,\ 1/1.414214) = (0.707107,\ 0.707107)\\
u\,\xhat/b &= 0.707107 \times (0.707107,\ 0.707107) = (0.5,\ 0.5)\\
\vect{x}/a - u\,\xhat/b &= (1 - 0.5,\ 0 - 0.5) = (0.5,\ -0.5)\\
\partial u/\partial\xhat &= (0.5,\ -0.5)/1.414214 = (0.353553,\ -0.353553)
\end{align*}
\textbf{The arccos factor.} $1 - u^2 = 1 - 0.5 = 0.5$; $\sqrt{0.5} = 0.707107$; $-1/0.707107 =
-1.414214$.

\textbf{Multiply.}
\[
\frac{\partial\theta}{\partial\xhat} = -1.414214 \times (0.353553,\ -0.353553) = (-0.5,\ +0.5).
\]
\textbf{Perpendicular to $\xhat$:} $(-0.5)(1) + (0.5)(1) = 0$.

\textbf{Length:} $\sqrt{0.25 + 0.25} = \sqrt{0.5} = 0.707107$, and $1/\norm{\xhat} = 1/1.414214
= 0.707107$. Equal.

Gradient descent moves $\xhat$ against the gradient, by $(+0.5, -0.5)$ times the step size:
it raises band 1 and lowers band 2, rotating $(1, 1)$ towards $(1, 0)$, as it should.
\end{solution}

\begin{solution}{ex:losses:3}
The derivative of arccos is $-1/\sqrt{1 - u^2}$, which grows without bound as $u \to 1$: $-22.4$
at $u = 0.999$, $-70.7$ at $0.9999$, $-\infty$ at 1. A well-trained model's pixels have angles of
0.02--0.05 rad, that is $u \ge 0.9987$, so the better the model, the larger this factor.

The true gradient does not blow up because the other factor, $\partial u/\partial\xhat$, has
length $\sqrt{1 - u^2}/\norm{\xhat}$, which shrinks at exactly the same rate. Their product has
length $1/\norm{\xhat}$ at every nonzero angle.

What goes wrong is the computation. Automatic differentiation evaluates the two factors
separately in float32, where the gap below 1 is $2^{-24} \approx 6 \times 10^{-8}$, so for
angles below about $2.4 \times 10^{-4}$ rad the computed $u$ is 1 or a hair above it. Then
arccos returns NaN (for $u > 1$) or the factor is $-\infty$ while $\partial u/\partial\xhat$ is
0, and $\infty \times 0$ is NaN. A clamp of $u$ to $[-1, 1]$ turns some of these NaNs into a
silent zero gradient but not all of them.

\code{SAMLoss} avoids arccos and computes $\theta = \operatorname{atan2}(\sin\text{-term},
\cos\text{-term})$. The guard that makes the derivative finite is the $+\varepsilon$ inside the
square root of the sine term (guard 3): it keeps that term at least $\sqrt{10^{-8}} = 10^{-4}$, so
the square root's derivative is bounded. Guard 2 (clamping the dot product to $[-N, N]$) keeps
the quantity under the root non-negative when rounding makes the dot product exceed the
product of lengths.
\end{solution}

\begin{solution}{ex:losses:4}
\textbf{Epoch index 4.}
\begin{align*}
t/R &= 4/10 = 0.4\\
\min(1, 0.4) &= 0.4\\
\lambda(4) &= 0.5 \times 0.4 = 0.20
\end{align*}
\textbf{Epoch index 50.} $50/10 = 5$, $\min(1, 5) = 1$, $\lambda = 0.5 \times 1 = 0.5$.
\begin{align*}
L &= 0.02 + 0.5 \times 0.12\\
&= 0.02 + 0.06\\
&= 0.08
\end{align*}
The angle contributes $0.06$ against L1's $0.02$: $0.06/0.02 = 3$ times more, 75\% of the
total.
\end{solution}

\begin{solution}{ex:losses:5}
\textbf{At $c = 0.12$.} Differences $y_i - c$: $-0.02, -0.01, 0.00, +0.01, +0.78$.
\begin{align*}
S_1 &= 0.02 + 0.01 + 0.00 + 0.01 + 0.78 = 0.82\\
S_2 &= 0.0004 + 0.0001 + 0 + 0.0001 + 0.6084 = 0.6090
\end{align*}
\textbf{At $c = 0.272$.} Differences: $-0.172, -0.162, -0.152, -0.142, +0.628$.
\begin{align*}
S_1 &= 0.172 + 0.162 + 0.152 + 0.142 + 0.628 = 1.256\\
S_2 &= 0.029584 + 0.026244 + 0.023104 + 0.020164 + 0.394384 = 0.493480
\end{align*}
$S_1$ is smaller at $c = 0.12$ (the median): $0.82 < 1.256$. $S_2$ is smaller at $c = 0.272$
(the mean): $0.4935 < 0.6090$. Squared error moves the best prediction towards the glitch;
absolute error does not.
\end{solution}

\begin{solution}{ex:losses:6}
Each prediction-mask flag lives on a token that covers a 4 $\times$ 4 block of pixels, and the
mask must stay strictly 0 or 1 so that \code{loss_mask == 1} selects whole blocks.
Nearest-neighbour upsampling copies each flag to its 16 pixels. Bilinear upsampling averages
neighbouring flags, so pixels near the boundary between a hidden and a visible token get
fractional values such as 0.25, 0.5 or 0.75. The test \code{loss_mask == 1} rejects them
silently, so a lattice of pixels along every block edge is never graded, and the product with
the eroded mask would no longer be a clean logical AND.
\end{solution}
